\documentclass[conference]{IEEEtran}

\usepackage{booktabs}
\usepackage{cite}
\usepackage{url}

\title{CudaNav: A Reproducible End-to-End GPU Autonomy Stack}

\author{\IEEEauthorblockN{Anonymous Authors}}

\begin{document}
\maketitle

\begin{abstract}
GPU robotics systems often benchmark isolated kernels while leaving data
movement, middleware integration, and end-to-end behavior outside the
measurement boundary. We present CudaNav, an open-source stack connecting GPU
LiDAR odometry, rolling voxel mapping, a typed Euclidean signed distance
field, and model predictive path integral control through reusable CUDA cores
and ROS~2/Nav2 adapters. We separate causal closed-loop simulation from
recorded-data shadow execution and from physical cross-device reproduction.
On an NVIDIA GeForce GTX~1660~Ti, the native stack completes 30 continuous
alternating S-course traversals over 1,059.4 simulated seconds and 352.75~m
with zero collisions, 0.0035\% final odometry drift, and zero frame-deadline
misses. On a content-addressed MCD Ouster sequence, GPU odometry processes
1,190 frames with 0.815~m ATE RMSE; the complete
odometry--mapping--ESDF--MPPI shadow stack obtains 0.812~m ATE RMSE and
integrates every scan. A release-profile ROS~2 simulation completes 30/30
traversals with retained MCAP and video evidence. Every result is bound to
source, data, hardware identity, thresholds, and raw outputs. Reproduction on
a second distinct physical GPU model remains an optional extension and is not
claimed by the present single-GPU evaluation.
\end{abstract}

\section{Introduction}
Autonomous navigation is a causal pipeline. LiDAR odometry changes the pose
used by mapping, mapping changes the distance field used by control, and
control changes the next state and scan. Accelerating correspondence search
or rollout evaluation alone does not establish that this chain is coherent,
safe, or deadline compliant.

CudaNav joins four reusable GPU stages:
\begin{center}
\fbox{\parbox{0.91\columnwidth}{\centering
LiDAR or recorded PointCloud2 $\rightarrow$ GPU KISS-ICP-style odometry
$\rightarrow$ rolling voxel map $\rightarrow$ typed ESDF
$\rightarrow$ CUDA MPPI $\rightarrow$ command-driven plant or shadow output}}
\end{center}
The stack exposes the same CUDA sources through native runners, Python
bindings, and ROS~2 lifecycle/Nav2 components. Its evidence contracts prevent
recorded motion from being reported as closed-loop success.

Our contributions are:
\begin{enumerate}
  \item an all-GPU odometry, mapping, distance-field, and control path without
        a host-side algorithmic fallback;
  \item ROS~2/Nav2 adapters that reuse the tested CUDA cores;
  \item content-bound contracts distinguishing simulation, real-data shadow
        execution, and physical multi-GPU reproduction; and
  \item a consumer-GPU evaluation that retains partial and negative evidence
        instead of silently dropping failed gates.
\end{enumerate}

\section{Related Systems}
KISS-ICP shows that point-to-point ICP with carefully selected robust
components can provide accurate LiDAR odometry across sensors
\cite{vizzo2023kissicp}. CudaNav follows this minimal design philosophy while
using GPU voxel-hash nearest-neighbor correspondences.

MPPI performs sampling-based receding-horizon control efficiently on GPUs
\cite{williams2017mppi}. CudaNav retains rollout validity, all-colliding,
retreat, and deadline diagnostics rather than reporting solve time alone.
Navigation2 provides lifecycle-managed ROS~2 navigation and behavior-tree
orchestration \cite{macenski2020marathon2}; CudaNav integrates at its costmap
and controller interfaces while keeping native and middleware evidence
separate.

\section{System}
\subsection{GPU LiDAR Odometry}
The odometry core downsamples each scan, predicts motion, searches neighboring
voxel-hash cells, rejects outliers, and solves a robust point-to-point rigid
update. Persistent device buffers retain the local map. Diagnostics include
correspondences, inliers, nearest-neighbor time, frame time, reference error,
and hash-capacity failures.

\subsection{Rolling Mapping and ESDF}
Registered points are fused into a bounded rolling voxel representation. A
typed distance-field interface carries dimensions, resolution, origin,
storage, and device ownership. GPU Euclidean distance propagation and
inflation produce the controller field, while observed-voxel, occupied-cell,
and map-shift counts expose stale or empty maps.

\subsection{CUDA MPPI}
The controller samples trajectories, evaluates path, goal, collision,
clearance, curvature-speed, and path-angle terms, and updates commands by
cost-weighted aggregation. DiffDrive, Ackermann, and Omni adapters share the
same core. The native S-course alternates path direction without resetting the
plant, odometry, voxel map, or ESDF, so accumulated drift remains visible.

\subsection{ROS 2 Boundary}
Lifecycle components provide GPU KISS-ICP, voxel mapping, and ESDF; a typed
distance-field message connects the Nav2 voxel costmap layer and CUDA MPPI
controller plugin. Exact-commit Jazzy CI checks package compilation, plugin
loading, parameter validation, derived-path sidecars, and no-GPU contract
tests.

\section{Evidence Contracts}
The evaluation distinguishes three non-interchangeable modes:
\begin{table}[t]
\centering
\caption{Evidence boundaries.}
\label{tab:evidence}
\small
\begin{tabular}{lcc}
\toprule
Mode & Causal scans & Real sensor \\
\midrule
Native/ROS~2 simulation & yes & no \\
Recorded-data shadow & no & yes \\
Physical GPU matrix & run-dependent & run-dependent \\
\bottomrule
\end{tabular}
\end{table}
Each manifest records a full commit, clean state, source digest, GPU model and
UUID, driver, profile, thresholds, commands, and hashes of retained files.
Validators reopen raw JSON, CSV, bag metadata, and media.

The machine-readable ledger tracks nine claims:
\texttt{reproducibility\_contracts},
\texttt{real\_dataset\_materialization},
\texttt{real\_gpu\_odometry},
\texttt{real\_gpu\_core\_shadow},
\texttt{native\_gpu\_core\_closed\_loop},
\texttt{integrated\_gpu\_stack},
\texttt{closed\_loop\_autonomy},
\texttt{real\_data\_shadow}, and
\texttt{multi\_gpu\_reproduction}. The first eight are supported. The final
claim is partial because only one physical GPU model is currently retained.

\section{Evaluation}
\subsection{Continuous Native Closed Loop}
The deterministic S-course uses a bounded $11\times5$~m workspace, two
offset wall obstacles, a 0.24~m robot radius, 240 LiDAR rays at three heights,
and a 0.1~s control period. Ground truth generates LiDAR and scores the run;
the controller receives only estimated state and GPU-produced map data. The
release profile requires 30 alternating traversals, at least 600 simulated
seconds, zero collisions, final goal distance below 0.30~m, drift below 1\%,
and frame-deadline miss rate below 1\%.

\begin{table}[t]
\centering
\caption{Native continuous closed-loop release.}
\label{tab:native}
\small
\begin{tabular}{lr}
\toprule
Metric & Result \\
\midrule
Traversals & 30/30 \\
Simulated duration & 1,059.4 s \\
Ground-truth distance & 352.748 m \\
Final goal distance & 0.296 m \\
Collisions & 0 \\
KISS-ICP ATE RMSE & 0.0124 m \\
Final drift & 0.003493\% \\
MPPI solve p95 & 0.455 ms \\
Full frame p95 & 5.237 ms \\
Deadline misses & 0\% \\
\bottomrule
\end{tabular}
\end{table}
The command-effect distance equals ground-truth travel, establishing that
commands affect subsequent plant state and scans. A second run on the same GTX
1660~Ti also completes 30/30 traversals, but it is repeat evidence rather than
cross-model evidence.

\subsection{Real PointCloud2 Shadow Execution}
The MCD \texttt{ntu\_day\_02} Ouster OS1-128 sources, calibration, and ground
truth are content addressed and deterministically materialized into a
7,559,663,616-byte ROS~2 SQLite/CDR bag. Full-window timing admission identifies
nanoseconds as the only plausible point-time unit. After the declared
one-second warm-up, 1,190 scans spanning 118.902~s remain.

\begin{table}[t]
\centering
\caption{MCD release-profile results.}
\label{tab:mcd}
\small
\begin{tabular}{lrr}
\toprule
Metric & Odometry & Full shadow stack \\
\midrule
Frames & 1,190 & 1,190 \\
ATE RMSE & 0.815 m & 0.812 m \\
Final drift & 0.472\% & 0.471\% \\
Mean frame time & 327.148 ms & 232.294 ms \\
GPU NN p95 & 177.821 ms & 180.835 ms \\
ESDF p95 & -- & 1.068 ms \\
MPPI solve p95 & -- & 0.715 ms \\
\bottomrule
\end{tabular}
\end{table}
The full stack integrates all scans, reaches 1,791,642 observed voxels, and
executes 120 two-iteration MPPI shadow evaluations without an all-colliding
event. These commands do not alter the recorded sequence and therefore do not
prove real-world closed-loop navigation.

\subsection{ROS 2 Closed Loop and Public Shadow Bag}
The ROS~2 release-profile simulator completes 30/30 traversals over
1,325.5~s and 380.12~m with zero collisions, 0.00283\% drift, and a 0.241\%
deadline-miss rate. A retained 1.23~GB MCAP contains positive counts for
command, odometry, occupancy, and ESDF topics, and a bound video records the
mission.

The public Istanbul shadow release adds 793 MPPI diagnostic cycles over
79.025~s. All 790 point clouds in the control window pair with a command
within 200~ms. Solve p95 is 4.801~ms, mean rollout validity is 0.9907, and
front clearance remains at least 3.535~m. Five retreat cycles are retained
rather than hidden.

\section{Reproducibility and Readiness}
The native and real-data release commands, thresholds, source hashes, dataset
hashes, and raw-output hashes are included in the artifact ledger. The strict
multi-GPU workflow executes the 30-traversal node on two separately labeled
self-hosted runners, then requires two distinct UUIDs and two distinct model
names at one commit and source digest.

Current evidence contains one NVIDIA GeForce GTX~1660~Ti node. The
\texttt{multi\_gpu\_reproduction} claim remains partial and optional until a
second distinct physical model passes the same release profile.
Real-robot closed-loop navigation is also outside the current claim boundary.
Neither limitation is represented as a completed result.

\section{Conclusion}
CudaNav demonstrates an end-to-end GPU autonomy path with causal native and
ROS~2 simulation, real PointCloud2 shadow execution, and content-bound
evidence. The retained results support the integrated single-GPU stack. The
anonymous source, PDF, and content-bound submission artifact are reproducible
without claiming cross-model GPU validation.

\bibliographystyle{IEEEtran}
\bibliography{references}

\end{document}
